\documentclass[12pt,leqno]{article}
\usepackage{amsmath,amssymb,amsthm,colortbl,xcolor,nicefrac,setspace}
\usepackage[T1]{fontenc}
\usepackage[expert,seriftt,lucidasmallscale]{lucidabr}
\usepackage{mathrsfs,hyperref,color,graphicx}
\usepackage{longtable,booktabs,array}
\DeclareMathSymbol{\Dashv}{\mathrel}{symbols}{239}
\DeclareMathSymbol{\eqdef}{\mathrel}{symbols}{214}
\DeclareMathSymbol{\sentails}{\mathrel}{symbols}{240}
\DeclareSymbolFont{stmry}{U}{stmry}{m}{n}
\DeclareMathDelimiter{\llbracket}{\mathopen}{stmry}{"4A}{stmry}{"71}
\DeclareMathDelimiter{\rrbracket}{\mathclose}{stmry}{"4B}{stmry}{"79}
\setlength{\textwidth}{6.5in}
\setlength{\textheight}{9in}
\setlength{\oddsidemargin}{0mm}
\setlength{\evensidemargin}{0mm}
\setlength{\topmargin}{0mm}
\setlength{\headheight}{0mm}
\setlength{\headsep}{0mm}
\setlength{\topskip}{0mm}
\setlength{\hoffset}{0in}
\setlength{\voffset}{0in}
\newcommand{\AND}{\ensuremath{\mathbin{\&}}}
\newcommand{\OR}{\ensuremath{\vee}}
\newcommand{\IFF}{\ensuremath{\leftrightarrow}}
\newcommand{\IC}{\ensuremath{\rightarrow}}
\newcommand{\NOT}{\ensuremath{\text{$\thicksim$}}}
\newcommand{\sem}[1]{\ensuremath{\llbracket #1\rrbracket}}
\newcommand{\Pure}{\mathsf{Pure}}
\newcommand{\Fun}{\mathsf{Fun}}
\newcommand{\funp}{\mathsf{fun}'}
\newcommand{\PP}{\mathsf{PP}}
\newcommand{\QSS}{\mathsf{QSS}}
\newcommand{\QLN}{\mathsf{QLN}}
\newcommand{\Ltwo}{\mathsf{L2}}
\newcommand{\Inv}{\mathsf{Inv}}
\newcommand{\WI}{\mathsf{WI}}
\newcommand{\TU}{\mathsf{TU}}
\newcommand{\RS}{\mathsf{RS}}
\newcommand{\PC}{\mathsf{PC}}
\newcommand{\NC}{\mathsf{NC}}
\newcommand{\id}{\mathsf{id}}
\newcommand{\Pow}{\mathcal P}
\newcommand{\Bmin}{\mathcal B_0}
\newcommand{\BR}{\mathcal B_R}
\newcommand{\BQ}{\mathcal B_Q}
\newtheorem{proposition}{Proposition}
\onehalfspacing
\setlength{\emergencystretch}{2em}
\brokenpenalty=10000
\hypersetup{colorlinks=true,linkcolor=blue!50!black,urlcolor=blue!50!black,
  pdftitle={Goodman's Purity of Pure Project: Verified Results and Open Problems},
  pdfauthor={},pdfsubject={Verified results, corrections, and open contributor problems}}

\begin{document}
\begin{center}
{\Large\bfseries Goodman's Purity of Pure Project\par}
\vspace{4pt}
{\large Verification, Corrections, and Remaining Problems\par}
\vspace{6pt}
20 September 2026; background updated 26 September 2026
\end{center}

\noindent
This report describes an Isabelle formalization of Jeremy Goodman's
project notes on Purity of Pure in higher-order metaphysics. It presents
the verified object-language results and model constructions, explains
corrections and qualifications to the arguments in the notes, and identifies
the open problems. The formalization builds on the independently maintained
Bacon--Dorr development of higher-order logic and Classicism.

Goodman's consistency question remains open. We have verified substantial
object-language arguments and exact-model constructions, including the
failure of L2 for the complete closed-logical stock, the rebuilt M5 model,
and a corrected M5 collision argument. We have not constructed a model of
the full background with the required instance of Purity of Pure, or
derived a contradiction from that background. The remaining mathematical
obligations are explicitly retained as contributor projects.

The software, source correspondence, and detailed claim ledger are
in the \href{https://github.com/fitelson/hom-isabelle/tree/main/Applications/goodman-isabelle}{Goodman application folder},
with the background logic supplied by
the \href{https://github.com/fitelson/hom-isabelle}{Bacon--Dorr core repository}
\cite{book,classic}.
The application shares the main repository's visibility and access permissions.
The report gives the mathematics rather than duplicating
the repository's implementation guide. Source labels T1--T9 and M1--M7
refer to Goodman's July 2026 notes \cite{goodman}.

\section{The question and the relevant distinctions}

Let $t$ be proposition type and put $u=t\to t$. For each type $\sigma$,
$\Pure_\sigma$ and $\Fun_\sigma$ have type $\sigma\to t$. The instance of
Purity of Pure used in the central arguments is
\[
 \PP_u := \Pure_{u\to t}(\Pure_u).
\]
It is one type-indexed instance, not the assertion of PP at every type.
Purity of Fun is a different principle and is not assumed in the question.

For pure operators $X,Y:u$, define
\[
 \funp(p) := \forall X\,Y\,
 ((\Pure(X)\AND\Pure(Y)\AND Xp=Yp)\IC X=Y).
\]
Thus evaluation at $p$ separates pure unary operators. This is not a
definition of $\Fun(p)$. Let $G$ be the pure operators with a pure
two-sided inverse, and define
\[
 X\approx Y := \exists Z\in G\;(X=Y\circ Z).
\]
Composition is on the input side. Identity at each type is the
object-language identity, not merely agreement of truth values.

For clarity, let $\Bmin$ denote the added logical-purity and
application-closure schemas over the represented CEV+ calculus. Logical
purity says that every closed term containing only logical vocabulary
denotes a pure entity; application closure says that applying a pure
function to a pure argument yields something pure. Let $\BR$ additionally
include exactly one fundamental proposition, no fundamentals at other
types, and zeroary and unary Recombination. Let $\BQ$ add the corresponding
Exhaustion principles. These distinguish the Recombination-only variant
from the full QLN version of the question. In the latter, the target is
consistency of $\BQ+\PP_u$, without Purity of Fun. Results for the weaker
variant must be identified separately.

Only zeroary and unary QLN are nonvacuous when there is exactly one
fundamental entity and fundamental arguments must be pairwise distinct.
This reduction of higher-arity instances is a mathematical observation
about the stipulated setting, not an additional cataloged theorem.

The implementation distinguishes ordinary consequence from an axiom
extension in which theorem-level rules operate above the added axioms.
In particular, a temporary assumption $\funp(r)$ cannot simply be promoted
to a permanent axiom. Nor may a biconditional established under a
temporary hypothesis automatically be converted to an identity by the
theorem-level Rule of Equivalence. The M5 correction below depends on this
distinction.

Proposition-valued Modalized Functionality, for $X,Y:\sigma\to t$, is
derived in the represented CEV development. Its arbitrary-result-type
extension is verified in the exact semantics but must not thereby be
called redundant over bare CEV. Likewise, model validity and
object-language derivability remain different claims.

\section{Scope of the formal verification}

The development establishes whole-proof forward preservation into the
named axiom-extension calculus, with its typing, signature, and closed-axiom
conditions explicit. This is not an unqualified two-way identification of
every presentation. The exact-model development proves the requisite
interpretation and global soundness conditions, including preservation
of the complete closed-logical denotational stock. Truth preservation of
sentences alone would not establish that stock correspondence.

All 303 supplied theory files are selected in the verified build. There are
1,708 integration catalog entries in 87 audit catalogs, with nine historical
replay entries counted separately. The audits check for oracle dependencies,
residual kernel hypotheses, and unresolved unification pairs. These counts
are not counts of independent mathematical results, and a clean audit does
not erase explicit theorem premises. Source fidelity requires the separate
comparison of those premises and definitions with the publications.

The semantic results use Isabelle/HOL with HOL--ZF. Their consistency
consequences are relative to that foundational setting. The exact appendix
construction is Bacon's finite-natural-word instance, with its constrained
function domains and the chosen singleton individual domain, not a theorem
about every monoid or every individual carrier. Its function domains satisfy
Bacon's stated action-compatibility condition.

\paragraph{The completed background metatheory.}
As of 26 September 2026, the core also verifies soundness of Bacon's
full-type Classicism for the nontrivial modal models and completeness
for signatures with countably many declared constants at each type.
The latter covers open formulas and infinite premise sets: derivability
is equivalent to consequence at the root, and consistency to
satisfiability. The results retain a rich variable stock, the minimal
logical vocabulary, inhabited domains, and a false proposition at every
world, with the documented future-restricted implication clause.
Soundness itself does not require countability.

These results strengthen the available background metatheory without
settling Goodman's question. Ordinary root consequence must still be
distinguished from the axiom extension used here, where theorem-level
rules operate above the added axioms. The existing closure result permits
the new completeness theorem to be applied to the full deductive closure
of the added stock; it does not establish that closure's consistency or
replace the global-validity conditions of the model results below.
No claim about $\Pure$ or $\Fun$ is strengthened merely by this update.

\subsection{The object-language results}

The following table summarizes the main numbered claims. The claim
ledger supplies their exact endpoints and axiom packages.

\setlength{\tabcolsep}{4pt}
\renewcommand{\arraystretch}{1.08}
\begin{longtable}{@{}>{\raggedright\arraybackslash}p{0.65in}
  >{\raggedright\arraybackslash}p{5.65in}@{}}
\toprule
Item & Verified content and qualification\\
\midrule
\endhead
T1 & With zeroary Exhaustion, pure propositions are truth or falsity.
The associated collapse of biconditional operators is verified.\\
T2a--c & Closure of $\funp$ under pure reversibles, nontriviality of
$\funp$ propositions, and attainment of each pure proposition are verified.
The displayed $\funp$ antecedents remain explicit.\\
T2d--f & Possible purity, false-but-possible noncontingency, and the six
pairwise-distinct propositions are verified in their stated scopes.
Possible purity has a proof not requiring Persistence; the six-distinctness
claim retains $\funp(r)$.\\
T3 & The advertised modal argument reaches possible identity where identity
is needed. The repaired heredity claims with zeroary Exhaustion or the
specified pure-identity rigidity principle are verified, retaining the
stated QSS and purity/persistence conditions. A defect in the
modal argument is not a countermodel to the whole CEV+ theory.\\
T4 & For the stated pure $C:t\to(t\to t)$, the higher-type failure of
$\funp$ at $C(r)$ is verified. The represented result needs logical purity
and application closure, not PP or $\funp(r)$.\\
T5 & The proliferation conclusion is verified as a conditional theorem
with the required $\funp(r)$ antecedent. It is not an existence or
consistency proof for that antecedent.\\
T6 & All four contradiction routes are verified with their L2 and
classification assumptions retained. Native formulations of the extra
principles and the repaired fundamental-witness routes are included.\\
T7 & T7a's absorption result is verified with its L2 assumptions.
T7b remains underspecified: the kind-level term, binders, and fixed/shifted
relations are not uniquely supplied by the notes.\\
T8 & The five base kinds, kind uniqueness under L2, and the stated
31-element growth bounds are verified. This is not by itself an unbounded
iteration proving infinitely many kinds.\\
T9 & The abstract conditional counting theorem and a substantive
exact-carrier instantiation, including conditional infinitude, are verified.
The arbitrary-model connection remains open; see Section~\ref{sec:attacks}.\\
\bottomrule
\end{longtable}

There are two further completed results concerning the notes' preliminary
claims. First, for each type $\sigma$ and each finite length $n$, pure
parameters $a_1,\ldots,a_n$ determine a pure predicate
\[
 P_n=\lambda x.\,(x=a_1\OR\cdots\OR x=a_n),
 \qquad P_0=\lambda x.\bot,
\]
with the expected material membership biconditional. The proof is an
explicit induction, using the closed logical function
$\lambda a.\lambda Q.\lambda x.(x=a\OR Qx)$ and application closure.
The empty case is separately verified. Neither PP nor Exhaustion is needed.
This proves finite pure comprehension, not comprehension for arbitrary
infinite external subsets. Its weaker axiom package inherits relative
consistency from the verified stronger no-PP QLN background.

Second, for $B_a=\lambda p.(p\IFF a)$, the identity
$B_a\circ B_a=\id$ is verified with no added axioms. The companion
purity and group-membership implications from $\Pure(a)$ are separately
packaged over the PP core. That is the package used in those proofs, not
a claim that PP is necessary.

The witness $R_+=\lambda p.(p\AND\funp(p))$ has a native proof of
purity, instantiation, and rigid specification. Consequently TU yields RS
over $\Bmin+\PP_u+\mathrm{Exhaustion}_0+\exists p\,\funp(p)+\TU$.
The contradiction between the notes' assertion that TU entails RS and their
assertion of incomparability remains a source question. Our theorem does
not establish that Goodman intended one assertion only in the weaker
setting without Exhaustion.

\section{The L2 principles and the open consistency question}
\label{sec:l2}

The notes have one L-principle in weak and strong forms, not a separate L1.
Weak L2 says
\[
 \Pure(X)\AND\Pure(Y)\AND\funp(p)\AND\funp(q)\AND Xp=Yq
 \quad\IC\quad X\approx Y.
\]
Strong L2 requires the same $Z\in G$ also to satisfy $q=Zp$.
The verified T6 consequences, writing the minimal added schemas explicitly,
are
\[
\begin{array}{l}
 \Bmin+\PP_u+\exists\funp+\Ltwo+\Inv\ \vdash\ \bot,\\
 \Bmin+\PP_u+\exists\funp+\Ltwo+\WI\ \vdash\ \bot,\\
 \Bmin+\PP_u+\exists\funp+\Ltwo+\TU\ \vdash\ \bot,\\
 \Bmin+\PP_u+\text{strong L2}+\RS\ \vdash\ \bot.
\end{array}
\]
Here $\exists\funp$ abbreviates $\exists p\,\funp(p)$. These are
axiom-extension derivations, not unrestricted uses of a deduction theorem.

The repaired route from fundamentality derives $\exists\funp$ using
Recombination, zeroary Exhaustion, and unique fundamentality. Recombination
supplies the unboxed identity information, but the available QSS argument
needs the further modal step supplied by zeroary Exhaustion. We do not
infer nonderivability of QSS from Recombination alone merely from that gap.
The T6 derivations that assume $\exists\funp$ directly require no QLN axiom.

For a model sound for the relevant axiom extension and globally validating
the repaired background with PP, T6 excludes simultaneous global validity
of L2 and TU. This is not automatically the claim that one of those formulas
is false at the distinguished root. Nor does failure of one make the
remaining axioms consistent.

\subsection{L2 fails for Bacon's complete closed-logical stock}

In the prefix presentation of the finite-natural-word frame, put
\[
 Z(P)=\{w:(\exists n\;w^\frown\langle n\rangle\in P)
             \AND(\exists n\;w^\frown\langle n\rangle\notin P)\}.
\]
Word reversal relates this presentation to Bacon's division action.
Isabelle identifies $Z$ with the denotation of a closed logical term in
the exact carriers. It proves that $Z$ is surjective, is not injective,
and is right-cancellative for composition in the complete closed-logical
stock. The definition tests variation among all immediate successors of a
world, not just two distinguished children.

For example, taking
$P_S=\{w^\frown\langle0\rangle:w\in S\}$ gives $Z(P_S)=S$.
Also $Z(P)=Z(\NOT P)$, so $Z$ is not reversible. Composition closure and
right-cancellation ensure that if $p$ separates the pure operators, so
does $Zp$. A verified generic construction supplies such a $p$. But
\[
 \id(Zp)=Z(p).
\]
L2 would make $\id\approx Z$. Since the witnessing input transformation
is reversible, this would make $Z$ reversible, a contradiction. The
formalization connects this calculation to actual root falsity of the
quantified L2 formula in the specified no-PP interpretation. Strong L2
fails there as well.

This answers Goodman's first model-calibration question for the complete
fixed logical stock. It does not show that L2 fails in every enlargement
of Pure, or that PP cannot imply L2. Enlarging Pure changes both L2 and
$\funp$. Historical bounded Vampire searches for L2 and sufficient
intermediate claims returned no proof in the reported runs. Those failures
are not underivability results, were not rerun for this report, and add no
Isabelle certificates to the present verification.

\section{The exact model claims and their repairs}

\paragraph{M1.}
At proposition type, the classifier of the complete logical stock is
the noncontingency operator. Bottom-type PP is verified. At the next
relevant type, the classifier exists in the exact domain and PP is
equivalent to its membership in the next pure stock. Its unconditional
nonmembership has not been proved. This part of M1 remains open.

Bacon's footnote-59 diagonal is separately verified, both on the exact
carriers and through an arbitrary-carrier compositional interpretation
with the explicit CEV+ soundness conditions. The obstruction retains
Purity of Fun, QSS, and the fundamental witness with the stipulated
uniqueness. A bare H model does not automatically provide that soundness
interface. Since Purity of Fun is absent from Goodman's question, this
obstruction is not a negative answer to it.

\paragraph{M2.}
For each set $T$ of propositions, the classifier
$G_T(P)=\{i:i\cdot P\in T\}$ is an invariant unary value, and the
construction gives a bijection with the invariant unary values. Cantor's
theorem and explicit collisions show why interpreting all invariants as
pure is incompatible with QSS at a fundamental proposition. This rejects
that interpretation of purity, not logical purity itself.

\paragraph{M3, M4, and M6.}
Freeness relative to the complete logical stock, both extreme views of a
$\funp$ proposition, and the product-meagerness result are instantiated
on the exact carriers. M4 has an actual suitable branch-lift witness
for each exact-stock $\funp$ input. M6 has actual separation and
nonindependence results, including a strict-inclusion pair whose relation
is preserved by substitutions. The required premises concerning the
complete logical stock are instantiated in these exact-model results.

The condition that every view of one fixed $r$ be $\funp$ is proved
impossible and is not assumed in the M4/M6 proofs. It must not
be confused with the correctly guarded modal QSS principle. The branch
in M4 must be chosen suitably; the arbitrary missing target in the
countability part of M6 is not asserted to be $\funp$.

\paragraph{M5's rebuilt model.}
The rebuilt interpretation uses Bacon's exact carriers.
Its expanded pure stock includes denotations of all closed terms over
the additional operator, including abstraction and higher-order
quantification, not merely a finite list of applications. The rebuilt
fundamental proposition supports the explicitly verified no-PP QLN
background and QSS. The displayed exotic operator is pure, self-inverse,
and neither uniformly truth-preserving nor uniformly truth-flipping.
TU and Inv have root failures; WI nonderivability is established by its
theorem-level relation to TU, not an unproved root implication.

Two source qualifications matter. For the displayed no-echo argument,
the claim that the only prefixes are the empty word and the word itself
uses a word of length one. Further, there are only countably many
world-indexed displayed pairs, so countability alone cannot select one
outside a countable orbit. A uniform diagonal pair construction supplies
the needed pre-rebuild obstruction. The later rebuilt interpretation
changes the fundamental proposition rather than leaving an old witness
fixed while enlarging Pure. No PP model is obtained.

\subsection{The corrected M5 collision}

The source also proposes that $F(p)=(p\IFF\NC(p))$, where
$\NC(p)=\Box p\OR\Box\NOT p$, has equal values at $\NC(r)$ and
$\top$ whenever $\funp(r)$. That local assertion is false in the
verified exact model. Adding $\funp(r)$ as a permanent axiom does not
justify this local claim: the resulting stock is already inconsistent.
Nontriviality gives $r\neq\top$ as
a theorem and hence its necessitation, whereas attainment gives
$\Diamond(r=\top)$. These conflict without PP. A local assumption or
the existential sentence $\exists\funp$ is a different matter.

The repaired input is $q=\NC(\Box r)$. The following is an actual
object-language derivation, not merely a semantic collision calculation.

\begin{proposition}
Over $\Bmin$, CEV+ proves
\[
 \funp(r)\IC(q\neq\top\AND Fq=F\top),
 \qquad
 \funp(r)\IC\NOT\mathrm{Reversible}(F).
\]
\end{proposition}

\begin{proof}
Write $q=\Box\Box r\OR\Box\NOT\Box r$. S4 proves $q\IC\Box q$.
It also proves $\NOT\Box\NOT q$: axiom 4 gives
$\NOT q\IC\NOT\Box r$, so normality gives
$\Box\NOT q\IC\Box\NOT\Box r$, hence $q$. Axiom T simultaneously
gives $\Box\NOT q\IC\NOT q$. Thus $q\IFF\NC(q)$ is a theorem.
Both $Fq$ and $F\top$ are therefore theorems, and theorem-level
Equivalence gives $Fq=F\top$ with no added axioms.

Now assume $\funp(r)$ temporarily. Nontriviality gives $\NOT\Box r$;
attainment at pure truth gives $\Diamond\Box r$. Hence both disjuncts
of $q$ are false. By T, $\NOT q$ implies $\NOT\Box q$, that is,
$q\neq\top$. Combine this with the previously established output
identity and discharge the temporary assumption. Reversibility implies
injectivity, yielding the second conclusion.
\end{proof}

The independent audit checks that this proof does not depend on collision
theorems using a fixed $\funp$ axiom or on the resulting collapse. The separate semantic
collision at the proposition true everywhere except the root remains
valid, but is not substituted for this proof. No general collision method
can refute reversibility merely because a candidate fails TU: the rebuilt
M5 model has a nonuniform involution with no collision. A PP-specific
argument would require further premises.

\paragraph{M7.}
The complete-stock diagonal is verified on the exact carriers for
every typed input $R$: a proposition exists that is neither pure nor the
value of a pure unary operator at $R$. Reachability by all invariant
operators is equivalent to injectivity of $i\mapsto i\cdot R$.
For every exact-stock $\funp$ input that orbit map is noninjective, with
an explicit unreachable singleton; this includes the verified generic
fundamental seed. Identifying a particular intended Theorem-10.1 gluing
with such a seed is still a separate contributor problem. No claim about
every arbitrary gluing is made.

\section{Bacon's appendix metatheory}

Proposition 8 and Theorem 10.1 here refer to the appendix of
\emph{Logical Combinatorialism} \cite{bacon}, not the introductory book.
The exact recursion's all-type closure, action laws, and surjectivity are
verified for the represented construction. Native global soundness and
the exact complete-stock correspondence are also verified. The supplied
QLN interpretation is the specified unique-fundamental-proposition
construction; it must not be silently identified with every gluing of
constant assignments.

The name-parametric Theorem 10.1 provides, for each countable family
$(A_n)_{n\in\mathbb N}$ of suitably typed interpretations over an
arbitrary type of constant names, one glued interpretation $C$ such that
\[
 \langle n\rangle\cdot\sem{M}_C=\sem{M}_{A_n}
\]
for every closed named term $M$ of the $t$-generated fragment. The same
$C$ works for all such $M$. The proof codes each term's finitely many
constant names into strings; it does not inject the entire signature
into strings or infer that the signature is countable. On string names
the polymorphic evaluator agrees with the string evaluator. The family of branches
remains countable. Open-term transport and types involving $e$ are not
included in this extension.

The enumeration and frame-representation proofs are selected,
audited, and transferred to named sentences. For a fixed string signature
and a closed sentence $M$ in the fragment, the one complete interpretation
$C$ satisfies
\[
\begin{aligned}
 M\text{ is frame-satisfiable}
     &\IFF C,\langle\rangle\models\Diamond M,\\
 M\text{ is frame-valid}
     &\IFF C,\langle\rangle\models\Box M.
\end{aligned}
\]
Frame satisfiability means root truth under some typed interpretation of
the constants on the fixed frame; validity means root truth under every
such interpretation. The modal forms retain the rich-variable and typed-
assignment conditions. This is not syntactic H completeness or an effective
decision procedure. The new arbitrary-name gluing theorem does not
automatically generalize the countable sentence-enumeration theorem.

\section{Goodman's suggested attacks}
\label{sec:attacks}

\paragraph{1. Calibrate L2.}
The fixed complete closed-logical stock has been calibrated negatively,
as explained in Section~\ref{sec:l2}. Whether PP forces L2 in the intended
full theory remains open.

\paragraph{2. Cap or characterize $G$.}
The classification principles and their T6 uses are verified, as are the
M5 counterexamples without PP. The granularity analysis supplies a precise
sufficient condition: a surjective proposition operator that preserves
sameness of truth value is truth-preserving or truth-flipping, assuming
there are both true and false propositions. Reversibility alone does not
give this condition.

With unary QLN, logical purity and application closure, the locally
guarded agreement test for pure $Z$ and fundamental $r$ gives
\[
 \NC(Zr\IFF r)
 \quad\IFF\quad
 \bigl(\forall p\,(Zp\IFF p)\bigr)
 \OR \bigl(\forall p\,\NOT(Zp\IFF p)\bigr).
\]
This is an equivalence to the corresponding truth-uniformity condition,
not a weaker principle already derived from PP. Unary Exhaustion is part
of the reduction. The TU-to-RS result above also retains its stated
Exhaustion and existence assumptions.

\paragraph{3. Exclude finitely many kinds under PC.}
Let $K$ be the kinds and $P$ the pure unary values. The checked exact-carrier
interpretation, with its global purity axioms, root L2 and $\exists\funp$,
and full external unary PC, gives
\[
 |\Pow(K)|\leq |P|\leq |K\times G|.
\]
The additional verified separation argument excludes finite $K$: if
$|K|=n$, classifiers selecting subsets of sizes $0,\ldots,n$ occupy
different kinds, since the induced reversible action preserves those
cardinalities. Thus $K$ is infinite and $2^{|K|}\leq|G|$.

This is real progress beyond the notes' finite-or-exponentially-large
dichotomy, but under explicit semantic hypotheses. Full external PC
ranges over every external subset of the relevant pure values; it is not
silently supplied by an internal comprehension formula. The transfer to
arbitrary suitable Henkin models, including the necessary quotient and
selector constructions, remains open. No bound $|G|\leq|K|$ has been
derived from Goodman's central theory. Such a bound would conflict with
the displayed inequality, but the cardinal argument does not supply it.

\paragraph{4. Prove consistency.}
No model of the full target has been verified. The exact no-PP models
show that failures of L2 or TU are possible in specified interpretations;
they do not discharge PP's self-classification requirement. The verified
finite-support theorem says that inconsistency of the axiom extension
would already be witnessed by a finite set of added axioms. It is a
proof-theoretic finiteness result, not a construction of models for every
finite set or a semantic compactness proof. Exploratory models and
unfinished search programs are not included among this report's results.

Goodman's PC also must not be conflated with Bacon's stronger identity
premise in the Melianism discussion. Coextension with a pure property
does not by itself identify that property with $\Pure_\sigma$.
The relevant infinite-disjunction argument is in Bacon's main text
with notes 57--58; note 60 points to the appendix. We retain the
classifier-extension result without claiming an independently verified
general complete-Boolean-algebra or unconditional classifier-nonpurity
theorem.

\section{What remains for contributors}

The verified development includes the native TU--RS transfer, finite PC at
every type, biconditional-operator algebra, named frame representation and
necessity, and arbitrary-name closed-term Theorem 10.1. These results do not
finish the entire formalization. The following four projects remain open.

\begin{enumerate}
\item Resolve unconditional higher-type classifier nonmembership in M1.
The existing equivalence with PP does not decide either side.
\item Supply the arbitrary-model T9 construction, retaining the distinction
between full external PC and internally represented subsets.
\item Specify the intended glued fundamental proposition and prove the
freeness or $\funp$ property needed for its M3/M7 application.
\item Obtain precise formulations of T7b and the multiple-fundamental M4
proposal, and clarify the conflicting TU/RS source passages.
\end{enumerate}

The repository gives bounded starting points for these projects. Further
algebra, open-term transport, and broader model constructions are separate
extensions, not accomplishments implied by the completed tasks. In
particular, the verified scope does not encompass every determinate source
claim, and not every theory file has received an exhaustive fidelity audit.

\section{Logical distinctions and verification safeguards}

The corrected results address the main dangers identified at the end of
Goodman's notes. Input negation and output negation are distinct; local
truth is not identity; a theorem of a model is not automatically a theorem
of the object calculus; substitutions form a monoid, not a group; and
$\funp$ expresses separation, not surjectivity of an orbit. The M5
countability error and the global-versus-local collision error are concrete
instances, not merely hypothetical warnings.

The present contribution is therefore a verified collection of arguments,
repairs, and precise limitations. It gives Goodman exact evidence about
which proposed routes work and where further mathematics is needed. It
does not favor either answer to the consistency question by itself.

\begin{thebibliography}{9}
\bibitem{goodman} Jeremy Goodman. \emph{Is Purity of Pure consistent with
Bacon's Logical Combinatorialism?} Project notes, July 2026.
\bibitem{bacon} Andrew Bacon. Logical Combinatorialism.
\emph{The Philosophical Review}, 2020. Appendix numbering follows the
source copy identified in the repository.
\bibitem{book} Andrew Bacon. \emph{A Philosophical Introduction to
Higher-Order Logics}. Routledge, 2024.
\bibitem{classic} Andrew Bacon and Cian Dorr. \emph{Classicism}.
Draft dated 1 July 2022; the core formalization's source numbering refers
to this draft.
\end{thebibliography}
\end{document}
